\cleardoublepage
\section{Prime and Composite Numbers}

Some numbers stand alone. Others are made of parts. Prime numbers have only two factors: $1$ and themselves. Composite numbers have more. This is not just a label—it is a clue to how numbers work.

$2, 3, 5, 7, 11, 13,$ and $17$ are prime. They cannot be broken down. But $15$ is composite: it splits into $3$ and $5$. Learning to recognize prime and composite numbers helps us factor numbers, simplify fractions, and understand the building blocks of math.

\subsection{Example: Prime and Composite Numbers}

As you explore numbers further, you'll discover that some numbers have only two factors, while others have many. A \textbf{prime number} is a number greater than $1$ that has exactly two distinct factors: $1$ and itself. For example, $2, 3, 5, 7, 11, 13, 17,$ and $19$ are all prime numbers under $20$. In contrast, a \textbf{composite number} has more than two factors. For instance, $15$ is composite because it can be divided by $1, 3, 5,$ and $15$.

Let's practice identifying these. Consider the numbers $14, 17, 21, 23,$ and $25$. Among these, $17$ and $23$ are prime because they have only two factors. The others—$14, 21,$ and $25$—are composite. If you look at all the composite numbers between $10$ and $20$, you'll find $10, 12, 14, 15, 16, 18,$ and $20$. It's also important to note that the number $1$ is neither prime nor composite, since it has only one factor: itself. Recognizing prime and composite numbers is a foundational skill for understanding more advanced math concepts like prime factorization and divisibility.

\newpage
\subsection{Practice Questions}

\begin{enumerate}
  \item \textbf{Short Question} \\
  Is 11 a prime or composite number?
  
  \textbf{Answer:} \textbf{Prime} — only two factors: 1 and 11.

  \item \textbf{Long Question} \\
  Identify all the prime numbers between 10 and 20. Explain how you know.
  
  \textbf{Answer:} \\
  \textbf{11, 13, 17, 19} — they each have exactly two factors: 1 and themselves.

  \item \textbf{Challenge Question} \\
  Why is the number 1 neither prime nor composite?
  
  \textbf{Answer:} \\
  It has only \textbf{one factor} (itself), not two or more. So it's \textbf{neither}.
\end{enumerate}

\newpage
\subsection{Hand-drawing 1}
This is where you will see a hand-drawn answer to the previous question in \textbf{English}.
\begin{figure}[htbp]
  \centering
  \includegraphics[width=\linewidth]{example-image} % TODO: Update caption for actual content
  \caption{A hand-drawn answer in English.}
\end{figure}

\newpage
\subsection{Hand-drawing 2}
This is where you will see a hand-drawn answer to the previous question in \textbf{Korean}.
\begin{figure}[htbp]
  \centering
  \includegraphics[width=\linewidth]{example-image} % TODO: Update caption for actual content
  \caption{A hand-drawn answer in Korean.}
\end{figure}

\cleardoublepage
\subsection*{Short Question (English)}
Is 11 a prime or composite number?

\fullpageimage{images/q1_5_1-eng.jpg}

\newpage
\subsection*{Short Question (Korean)}
Is 11 a prime or composite number?

\fullpageimage{images/q1_5_1-kor.jpg}

\newpage
\subsection*{Short Question (Answer)}
Is 11 a prime or composite number?

\textbf{Answer:} \textbf{Prime} — only two factors: 1 and 11.

\cleardoublepage
\subsection{Long Question (English)}
Identify all the prime numbers between 10 and 20. Explain how you know.

\fullpageimage{images/q1_5_2-eng.jpg}

\newpage
\subsection{Long Question (Korean)}
Identify all the prime numbers between 10 and 20. Explain how you know.

\fullpageimage{images/q1_5_2-kor.jpg}

\newpage
\subsection{Long Question (Answer)}
Identify all the prime numbers between 10 and 20. Explain how you know.

\textbf{Answer:}\\
\textbf{11, 13, 17, 19} — they each have exactly two factors: 1 and themselves.

\cleardoublepage
\subsection{Challenge Question (English)}
Why is the number 1 neither prime nor composite?

\fullpageimage{images/q1_5_3-eng.jpg}

\newpage
\subsection{Challenge Question (Korean)}
Why is the number 1 neither prime nor composite?

\fullpageimage{images/q1_5_3-kor.jpg}

\newpage
\subsection{Challenge Question (Answer)}
Why is the number 1 neither prime nor composite?

\textbf{Answer:}\\
It has only \textbf{one factor} (itself), not two or more. So it's \textbf{neither}.

